% COLING 2027 submission via ACL Rolling Review (ACL template, review mode, anonymous).
\documentclass[11pt]{article}
\usepackage[review]{acl}
\usepackage{times}
\usepackage{latexsym}
\usepackage[T1]{fontenc}
\usepackage[utf8]{inputenc}
\usepackage{microtype}
\usepackage{inconsolata}
\usepackage{amsmath,amssymb}
\usepackage{booktabs}
\usepackage{multirow}
\usepackage{xcolor}

\newcommand{\ci}[2]{\,[{#1},\,{#2}]}
\newcommand{\todo}[1]{\textcolor{red}{[TODO: #1]}}
\newcommand{\model}{RoleNet}
\newcommand{\pb}{PaperBest}

\title{Watch the Listener: Identity-Free Responder Modeling\\for Emotion Forecasting in Multi-Party Video}

\author{Anonymous submission}

\begin{document}
\maketitle

\begin{abstract}
Emotion forecasting asks how a person will feel in the next conversational turn, before that turn is observed.
In the Hi-EF benchmark the future responder is not identified at inference time, and existing models summarise
the preceding clips around whoever is speaking. We argue that the responder is usually already on screen,
listening, and that a forecaster should be built around that person. We present \model{}, which (i) selects a
prospective responder without any future information and (ii) represents the preceding clips as a set of
participant-by-turn evidence tokens read by a set encoder. On the locked Hi-EF test split, evaluated once under a
preregistered plan, \model{} improves unweighted average recall over the strongest published configuration by
3.84 points \ci{1.24}{6.76}; the preregistered primary contrast under logit-adjusted scoring was not confirmed,
which we trace to a test set with six examples of one class. Replacing our responder proxy by the true responder
identity changes nothing ($-0.04$ \ci{-1.50}{1.85}), whereas removing the listener's face or the preceding
context each costs about two points. Explicit temporal, relational or role structure in the encoder brings no
measurable gain. Remaining errors concentrate on turns where the responder does not mirror the speaker's emotion.
\end{abstract}

% ======================================================================================================
\section{Introduction}

Recognising an emotion that is displayed is a perception problem; forecasting an emotion that has not yet been
displayed is a problem about people. Affective forecasting has a long history in psychology
\citep{wilson2003affective}, and \citet{wang2026hief} recently turned a version of it into a benchmark:
given three consecutive clips of a TV conversation (I--III), predict the emotion of the person who speaks in clip IV.
Two properties make the task unusual. The target clip is never observed, and the identity of the future speaker
(the \emph{responder}, B) is not given.

Existing systems for this setting process each clip as a whole and fuse the clips over time
\citep{wang2026hief}. The clip III speaker (A) dominates such summaries, because the camera, the transcript and the
audio follow whoever talks. Yet the person who will answer is often visible already, listening to A.
Two regularities of conversation suggest that this person matters most. Responders tend to take over the emotion
of their interlocutor \citep[emotional contagion;][]{hatfield1994contagion}, and individuals tend to stay in their
own recent emotional state \citep[emotional inertia;][]{kuppens2010inertia}. Both pass through the responder: the
first is visible as a reaction while listening, the second as the responder's own earlier expressions.

We build the forecaster around the responder. The difficulty is that the responder is unknown. Our first
contribution is an \textbf{identity-free responder proxy}: within the current turn, the second most visible person
(after the speaker) is taken as the listener, and all evidence is organised around this listener, the speaker and
everyone else. No future information, speaker label or identity annotation is used. Our second contribution is a
\textbf{participant-centric evidence set}: each participant in each preceding turn contributes one token built from
their face frames, alongside one speech token and one scene token per turn. A small set encoder reads these tokens
jointly and a query token pools them into a forecast. We call the model \model{}.

We evaluate under a strict protocol. All design decisions were made with episode-level cross-validation on the
training and validation episodes; the locked test split was then read once under a preregistered plan
(Section~\ref{sec:protocol}). On test, \model{} improves unweighted average recall (UAR) over the strongest
configuration of \citet{wang2026hief}, retrained on the same split, by 3.84 points \ci{1.24}{6.76}, and wins
seven of eight test episodes. We report openly that the preregistered primary metric, logit-adjusted UAR, did not
confirm this contrast, and why.

Controlled experiments then explain the gain. The proxy loses nothing against the true responder identity. The
listener's face and the immediately preceding context are each worth about two UAR points, while the encoder needs
no explicit temporal order, relation biases or role labels. Finally, an analysis of the remaining errors shows that
forecasts succeed mostly when the responder mirrors the speaker and fail when the responder shifts to a different
emotion, and that the missing information lies in recognising conversational states rather than in combining them.

% ======================================================================================================
\section{Related Work}

\paragraph{Emotion recognition and forecasting in conversation.}
Emotion recognition in conversation labels utterances that are observed, typically with speaker-aware memory or
graph models \citep{hazarika2018icon,majumder2019dialoguernn,ghosal2019dialoguegcn} and multimodal transformers
\citep{tsai2019mult}; multi-party benchmarks such as MELD provide speaker names \citep{poria2019meld}.
Forecasting the sentiment or emotion of a turn that has not happened yet has been studied in text dialogue
\citep{wang2020sentiment}. Hi-EF \citep{wang2026hief} brings the task to multimodal video with an unidentified
responder, which removes the speaker-tracking machinery that recognition models rely on.

\paragraph{Responders and listeners.}
DialogueRNN keeps a state per party, updated by utterances \citep{majumder2019dialoguernn}; it assumes that party
identities are known. We instead select the prospective responder from what is visible in the current turn, and
show that knowing the true responder adds nothing on Hi-EF.

\paragraph{Set models and interaction analysis.}
Attention over an unordered collection of elements with learned pooling seeds is formalised by the Set Transformer
\citep{lee2019set}; see also \citet{zaheer2017deep}. We use this view to describe \model{} and to test which structure
it needs. To measure whether a trained model uses cross-group interactions we follow EMAP \citep{hessel2020emap}.

% ======================================================================================================
\section{Task and Data}
\label{sec:task}

An MCIS (multilayered-contextual interaction sample) in Hi-EF \citep{wang2026hief} consists of four consecutive
clips from a TV series, each clip one subtitle utterance. Clips I and II are context, clip III is a turn of speaker
A, and clip IV is a turn of a different person B. The target is B's emotion in clip IV, one of seven classes
(angry, disgust, fear, happy, neutral, sad, surprise). Only clips I--III are inputs.

\paragraph{Split.} The released annotations give 2{,}830 MCIS from 53 episodes. To avoid shared clips between
partitions we split by episode (37/8/8 episodes; 1{,}993/428/409 MCIS for train/validation/test). Our numbers are
therefore not comparable to those printed in \citet{wang2026hief}, who used a different split; we retrain their
strongest configuration on ours (Section~\ref{sec:protocol}). Classes are imbalanced: the 409 test MCIS contain
six examples of \emph{fear}.

\paragraph{Constraints.} Clip IV is never an input at inference. B's identity is not given. We do not assume that
only two people take part.

% ======================================================================================================
\section{Method}
\label{sec:method}

\subsection{Participants and the responder proxy}

Faces are detected in frames sampled at 4\,fps (at most 32 per clip) and described by ArcFace identity embeddings
\citep{deng2019arcface}. Within an MCIS, embeddings from clips I--III are clustered (average linkage, cosine
threshold 0.45) into identities. In clip III, the identity with the most frames is taken as the speaker A and the
identity with the second most frames as the \emph{listener} L, our proxy for the responder B. Every other identity
is a bystander O. If clip III shows only one person, no listener is assigned. The rule uses no clip IV, no
transcript speaker label and no identity annotation.

\subsection{The participant-centric evidence set}

\paragraph{Face tokens.} For role $r \in \{A, L, O\}$ and clip $k \in \{1,2,3\}$, let $F_{r,k}$ be the frames of the
faces of role $r$ in clip $k$ (at most 24). Each frame $f$ is described by $\phi(f) \in \mathbb{R}^{146}$: a
128-dimensional PCA of the HSEmotion embedding \citep{savchenko2022hsemotion} (fitted on training faces of the
current fold), its 8 expression probabilities and valence/arousal, head pose, face box position and size, mouth
opening and relative time in the clip. Frames are pooled by attention,
\begin{equation}
u_f = \mathrm{MLP}(\phi(f)),\quad
\alpha_f = \operatorname{softmax}_{f \in F_{r,k}}\!\big(w^\top u_f\big),\quad
\bar u_{r,k} = \textstyle\sum_f \alpha_f u_f ,
\end{equation}
and the token is $x_{r,k} = \bar u_{r,k} + e^{\text{role}}_r + e^{\text{clip}}_k$, with a learned placeholder
$a_{r,k}$ replacing $\bar u_{r,k}$ when $F_{r,k}$ is empty.

\paragraph{Speech and scene tokens.} For each clip $k$, a speech token combines the CLIP text embedding of the
transcript \citep{radford2021clip}, the AudioCLIP sound-event vector \citep{guzhov2022audioclip} gated by an
audio-found flag, and a 9-dimensional turn descriptor $v_k$ (mouth--audio synchrony per role, ECAPA
\citep{desplanques2020ecapa} voice similarity to clip III, and the number of identities):
$s_k = W_t\,t_k + (W_a\,a_k)\,\delta_k + W_v\,v_k + e^{\text{speech}} + e^{\text{clip}}_k$.
A scene token $g_k = W_s[\bar c^{\text{frame}}_k;\bar c^{\text{face}}_k] + e^{\text{scene}} + e^{\text{clip}}_k$
uses mean CLIP embeddings of the whole frames and of all faces.

The evidence of an MCIS is the set of 15 typed elements
\begin{equation}
\mathcal{S} = \{x_{r,k}\}_{r \in \{A,L,O\},\,k \le 3} \;\cup\; \{s_k\}_{k \le 3} \;\cup\; \{g_k\}_{k \le 3}.
\end{equation}

\subsection{Set encoder and pooling}

A learned query $q$ is prepended and the set is processed by two pre-norm set attention blocks
\citep[SAB;][]{lee2019set}, i.e.\ transformer layers \citep{vaswani2017attention} without positional encoding:
\begin{equation}
Z = H + \mathrm{MHA}(\mathrm{LN}(H)),\qquad H' = Z + \mathrm{FFN}(\mathrm{LN}(Z)),
\end{equation}
with $H^{(0)} = [q; \mathcal{S}]$, $d = 128$, 4 heads and dropout 0.2. The forecast is
$p(y_B \mid \mathcal{S}) = \operatorname{softmax}(W\,\mathrm{LN}(h_q))$, read from the query position.
Because the query participates in self-attention it acts as a pooling seed in the sense of PMA \citep{lee2019set}.
The encoder is equivariant to permutations of the typed elements; order enters only through the clip embeddings,
which we show can be removed (Section~\ref{sec:structure}). The model has 0.71M parameters.

\subsection{Training}

The loss is cross-entropy on $y_B$ plus three auxiliary terms (weight 0.3 each): $y_B$ from the mean of the face
tokens, $y_B$ from the mean of the speech and scene tokens, and A's clip III emotion from A's clip III face token.
Modality dropout \citep{neverova2016moddrop} removes all speech and scene tokens with probability 0.3, or else all
face tokens with probability 0.15, which amounts to training on random subsets of $\mathcal{S}$. We use AdamW
(learning rate $3\cdot10^{-4}$, weight decay $10^{-2}$, batch 64, at most 80 epochs) and keep the checkpoint with the
best UAR on five held-out training episodes (patience 12). Predictions are averaged over seeds.

% ======================================================================================================
\section{Experimental Protocol}
\label{sec:protocol}

\paragraph{Baseline.} \pb{} is the strongest configuration of \citet{wang2026hief}: per-clip temporal transformers
over face and frame features, cross-attention fusion with text and audio, and LSTM plus transformer fusion across
clips. We retrain it on our split as re-implemented from the released code (AdamW, learning rate $10^{-4}$, batch
32, 50 epochs, checkpoint by selection UAR). As the best configuration reported for Hi-EF, it is our reference.

\paragraph{Development.} All design decisions use 5-fold cross-validation over the 45 training and validation
episodes (2{,}421 MCIS), with folds balanced by MCIS count. Unless stated otherwise we report seed ensembles of 10
seeds and 95\% intervals from a bootstrap that resamples seeds and episodes (2{,}000 draws).

\paragraph{Test.} The test split was read once for \model{} and \pb{}, under a plan committed before the run
(models, seeds, metrics and the order of confirmatory contrasts). Models were trained on the 45 development
episodes with five seeds. The preregistered primary metric was 7-class UAR after logit adjustment
\citep{menon2021logit} of the seed-averaged probabilities; plain argmax scoring, the metric of
\citet{wang2026hief}, was preregistered as descriptive. Intervals resample the 8 test episodes. Earlier,
unrelated model families had been evaluated on the same test split; Appendix~\ref{app:test} lists every access.

\paragraph{Metrics.} UAR (macro recall) and WAR (accuracy); for controlled comparisons on development data we also
report negative log-likelihood after temperature scaling \citep{guo2017calibration} on held-out training episodes.

% ======================================================================================================
\section{Results}

\subsection{Main result}

\begin{table}[t]
\centering\small
\begin{tabular}{lcc}
\toprule
Model (test, seed ensemble) & UAR & WAR \\
\midrule
\pb{} \citep{wang2026hief}, retrained & 21.40 & 32.52 \\
\model{} (ours) & \textbf{25.24} & \textbf{36.43} \\
\midrule
$\Delta$ plain (descriptive) & $+3.84$\ci{1.24}{6.76} & $+3.91$\ci{0.40}{7.79} \\
$\Delta$ logit-adjusted (primary) & $+3.77$\ci{-7.49}{10.12} & -- \\
$\Delta$ logit-adj., 6 classes & $+9.95$\ci{6.6}{13.6} & -- \\
\bottomrule
\end{tabular}
\caption{Locked Hi-EF test (409 MCIS, 8 episodes), read once. Intervals: 95\% bootstrap over test episodes.
\model{} wins 7 of 8 episodes (plain). The 6-class row excludes \emph{fear} (6 test examples).}
\label{tab:main}
\end{table}

Table~\ref{tab:main} shows the test results. Under plain scoring, \model{} improves UAR by 3.84 points and WAR by
3.91 points over \pb{}. The preregistered primary contrast, UAR after logit adjustment, has a similar point estimate
($+3.77$) but an interval that includes zero. The cause is the \emph{fear} class: with six test examples, one
correct \emph{fear} prediction is worth 2.38 UAR points, and logit adjustment pushes \pb{} to predict \emph{fear}
for 27.6\% of test MCIS (true rate 1.5\%), gaining two \emph{fear} hits while its \emph{angry} recall falls to
1.3\%. Without \emph{fear}, the adjusted contrast is $+9.95$ \ci{6.6}{13.6}. We therefore regard the test result as
consistent with, but not a confirmation of, the development-set evidence, and we flag the choice of an adjusted
7-class primary metric for a test set with six \emph{fear} examples as a design error of our preregistration.

\subsection{Does the forecaster need to know who responds?}
\label{sec:proxy}

To test the responder proxy we construct an oracle that replaces the listener by the true responder. B is found as
the dominant face of clip IV and matched to an identity of clips I--III (possible for 88\% of MCIS); this uses clip
IV and is available only for analysis. The proxy's listener equals the matched responder in 68\% of the MCIS where
the proxy assigns a listener.

\begin{table}[t]
\centering\small
\begin{tabular}{lrcc}
\toprule
Subset (development CV) & $n$ & Proxy & Oracle \\
\midrule
All MCIS & 2{,}421 & 26.24 & 26.20 \\
B seen in clips I--III & 2{,}107 & 26.98 & 26.90 \\
\quad proxy picked another person or none & 390 & 28.33 & 28.15 \\
B not seen in clips I--III & 314 & 21.20 & 21.29 \\
\bottomrule
\end{tabular}
\caption{UAR with the identity-free listener proxy vs.\ the true responder (oracle, uses clip IV). 10 seeds.
Oracle $-$ proxy: $-0.04$ \ci{-1.50}{1.85}.}
\label{tab:proxy}
\end{table}

Table~\ref{tab:proxy} shows no difference, also on the 390 MCIS where the proxy picks the wrong person.
Together with the next experiment, this indicates that the model uses \emph{what the visible faces show}, not
\emph{whose} faces they are: an identity-free selection of the evidence suffices.

\subsection{Which evidence matters?}
\label{sec:evidence}

\begin{table}[t]
\centering\small
\begin{tabular}{lcc}
\toprule
Variant (development CV, 10 seeds) & UAR & $\Delta$ vs.\ full \\
\midrule
\model{} (full) & 26.24 & -- \\
\quad without listener tokens & 24.20 & $-2.04$\ci{-3.46}{-0.21} \\
\quad clip III only (no I, II) & 24.11 & $-2.13$\ci{-4.08}{-0.20} \\
\quad without role labels & 25.80 & $-0.44$\ci{-2.16}{1.21} \\
\quad mean pooling instead of query & 26.42 & $+0.18$\ci{-1.62}{1.55} \\
\bottomrule
\end{tabular}
\caption{Evidence ablations, each retrained. Intervals resample seeds and episodes.}
\label{tab:evidence}
\end{table}

Two sources survive seed and episode noise (Table~\ref{tab:evidence}). Removing the listener's tokens costs 2.04
UAR (all five folds), and the cost is concentrated where a listener is visible in clip III ($-2.80$ on 1{,}451 MCIS
vs.\ $-0.54$ on 970). Removing clips I and II costs 2.13; clip II carries most of this context. Pooling all faces
without role labels is not significantly worse ($-0.44$): what matters is that the listener's face enters the
evidence set, not that it is labelled as such.

\subsection{Does the encoder need more structure?}
\label{sec:structure}

\begin{table}[t]
\centering\small
\begin{tabular}{lc}
\toprule
Change to \model{} (development CV) & Effect \\
\midrule
Separate additive role paths & UAR $-2.14$\ci{-3.79}{-0.29} \\
\quad + learned interaction term & UAR $-1.32$\ci{-2.80}{0.34} \\
Clips I/II unordered & NLL $+0.006$\ci{-0.010}{0.022} \\
Clips I--III unordered & NLL $+0.008$ \\
Relation biases (time, person, evidence) & NLL $+0.001$\ci{-0.014}{0.017} \\
Readout variants (4 kinds) & UAR within seed noise \\
Additive projection over L / history / event & NLL $+0.011$\ci{0.005}{0.016} \\
\bottomrule
\end{tabular}
\caption{Structural variants. Separate paths: 3 seeds. Order and relation biases: 10 seeds, temperature-scaled
NLL relative to \model{} (positive = worse). Readouts: mean pooling, listener token, read-only query, three queries.
Additive projection: group-level EMAP of the trained model, 5 seeds $\times$ 3 folds.}
\label{tab:structure}
\end{table}

We describe \model{} as a set encoder; Table~\ref{tab:structure} tests whether it needs more structure than that.
Forcing the evidence into separate role paths (listener, event, scene), each recurrent over clips and combined
additively, costs 2.14 UAR, and a learned interaction term recovers only part of it: the evidence must be processed
jointly. Making the encoder exactly invariant to the order of clips I and II, or of all three clips (verified
numerically, maximum logit change $1.4\cdot10^{-6}$), does not hurt. Factorial attention biases for temporal,
person and evidence-type relations are learned but bring no gain, nor do alternative readouts.
A group-level EMAP \citep{hessel2020emap} of the trained model shows that the encoder does use non-additive
interactions. We split the tokens into the listener (clips I--III), the history (other tokens of clips I--II) and
the current event (other tokens of clip III), and replace the model by its best additive approximation over these
groups. Held-out NLL rises by 0.011 \ci{0.005}{0.016}, and each group interacts with the rest ($+0.005$ each, all
intervals above zero). The interactions are modest, about 5\% of the logit variance, and they do not involve the
transcript: an additive split of text versus everything else loses nothing ($-0.002$ \ci{-0.004}{0.000}).

\subsection{Generalisation to MELD}

We rebuilt Hi-EF-style windows from MELD \citep{poria2019meld} (three turns, then a turn of a different speaker;
5{,}406/595/1{,}341 windows), without speaker names as input, and trained \model{} and \pb{} with three seeds
(train for fitting, dev for early stopping, test read once for the comparison). \model{} reaches 16.85 UAR vs.\
15.50 for \pb{} ($+1.35$ \ci{-1.02}{4.24}; chance 14.29). \pb{} predicts \emph{neutral} for most windows (recall
93\%), which gives it a higher WAR (46.09 vs.\ 39.30) and a lower NLL (1.543 vs.\ 1.695). The advantage on MELD is
therefore small and not confirmed.

% ======================================================================================================
\section{Analysis: Where Forecasts Fail}
\label{sec:analysis}

\paragraph{Mirroring and shifts.} B's emotion equals A's clip III emotion in 35.5\% of development MCIS (34.7\% in
our MELD windows). \model{} is correct on 56.9\% of these \emph{mirror} cases but on only 27.7\% of \emph{shifts},
where it still predicts A's emotion 29\% of the time. The listener's tokens help mostly on mirror cases. Whether a
shift will happen is partly predictable (AUC 0.70), but which emotion follows a shift is not: recall on shifts is
1.6\% for \emph{disgust}, 2.8\% for \emph{surprise} and 0\% for \emph{fear}.

\paragraph{The missing information is a recognition problem.} Two oracle analyses locate the headroom. With A's
true clip III emotion added to \model{}'s outputs, UAR rises from 24.73 to 27.51 (NLL $-0.036$
\ci{-0.049}{-0.021}); with A's emotion estimated by a recogniser from clip III, the gain disappears. On the 26\% of
MCIS where B's own earlier turn is labelled, simply copying B's earlier emotion gives 39.07 UAR against 27.98 for
\model{}, and adding \model{}'s outputs to that label adds nothing detectable. The useful signal is the
participants' current emotional states, which the fixed feature extractors recognise poorly; better ways of
combining the same tokens cannot recover it.

% ======================================================================================================
\section{Conclusion}

Forecasting a responder's emotion on Hi-EF is mostly a matter of looking at the right person. An identity-free
listener proxy and a participant-centric evidence set read by a small set encoder outperform the published
configuration on a locked test split, the proxy loses nothing against the true responder, and additional
structure in the encoder does not help. The remaining errors point to recognising the participants' current
emotional states, especially before a shift, as the next problem to solve.

% ======================================================================================================
\section*{Limitations}

\textbf{Primary metric not confirmed.} The preregistered primary contrast (logit-adjusted UAR) did not reach
significance; our main claim rests on plain UAR, which was preregistered as descriptive and is the benchmark
standard. \textbf{Test reuse.} The test split had been used before for unrelated model families
(Appendix~\ref{app:test}); \model{}'s design used development data only. \textbf{Small test set.} Eight test
episodes give wide intervals; one \emph{fear} example is worth 2.38 UAR points. \textbf{Exploratory development.}
The design of \model{} was informed by analyses on the same development episodes used for cross-validation.
\textbf{Responder oracle.} The true responder is approximated by the dominant face of clip IV matched by face
identity; matching errors would blur the proxy--oracle comparison. \textbf{Features.} Text and audio come from
general-purpose encoders (CLIP, AudioCLIP sound events), not from emotion- or prosody-specific models; conclusions
about which evidence matters are conditional on them. \textbf{Domain.} Hi-EF and MELD consist of scripted TV series
with shot/reaction-shot editing; the listener proxy relies on that editing style and may not transfer to fixed-camera
recordings. \textbf{MELD.} On MELD the gain is small and not confirmed, and the MELD test set was read twice: the
first run failed on three non-finite feature values, which were set to zero before the second run.

\section*{Ethical Considerations}

We use existing public benchmarks of scripted TV content. Face and voice embeddings are used only to group
observations within a short sample; no identity is linked to a name. Forecasting a person's emotion could be misused
for manipulation or surveillance; our models are research artefacts with modest accuracy (about 25 UAR on seven
classes) and should not be used to make decisions about individuals.

\bibliography{references}

\appendix

\section{Test-Split Accesses}
\label{app:test}

\todo{Table from SYNTHESIS \S1: G4 run 1 (recognize-then-transition and trajectory forecasters, B1), G4 re-run and
G5 (announced, execution not verified), G10 (this paper, single preregistered run). State that no RoleNet
hyper-parameter was chosen after any test access.}

\section{Hyper-parameters and Implementation}
\label{app:hparams}

\todo{Full RN dictionary, PCA fitting per fold, identity clustering thresholds (same-person cosine 0.45, dominant
fraction 0.25), voice threshold 0.35, frame caps, seeds, PaperBest settings, compute.}

\section{Full Ablation Tables}
\label{app:ablations}

\todo{G13 token ablations with the same-model noise distribution (220 three-seed ensembles, 25.55 $\pm$ 0.40), G33a
and G33b tables, G12 path shares.}

\end{document}
